\documentclass[aspectratio=169,10pt,english]{beamer}

\input{../../../_preambulo_beamer.tex}
\input{../../../_comandos_beamer.tex}

\selectlanguage{english}

\title{MA1001B - Statistical Modeling for Decision Making}
\subtitle{Lesson 41: Chi-Square Test of Independence}
\author{}
\institute{Tecnol\'ogico de Monterrey}
\date{}

\begin{document}

\begin{frame}[plain]
  \vspace{1.2cm}
  {\Large\bfseries MA1001B - Statistical Modeling for Decision Making\par}
  \vspace{0.7cm}
  {\large Lesson 41: Chi-Square Test of Independence\par}
  \vspace{0.7cm}
  {\color{TecRojo}\rule{\textwidth}{1.2pt}}
  \vspace{0.8cm}

  MA1001B Faculty\\[0.3cm]
  Term\\[0.3cm]
  Tecnol\'ogico de Monterrey
\end{frame}

\begin{frame}{The Two-Way Association Question}
  \begin{alertblock}{Guiding question}
    In one sample of $n=325$ visitors, is conversion independent of
    acquisition channel?
  \end{alertblock}

  \vspace{0.6em}
  By the end of this lesson, you can:
  \begin{itemize}
    \item form expected counts from the table margins;
    \item compute $\chi^2$ with $df=(r-1)(c-1)$;
    \item report a p-value and Cramer's V;
    \item refuse a causal reading of a significant association.
  \end{itemize}

  \vfill
  \scriptsize All visitor counts used today are synthetic. No real company data are used.
\end{frame}

\begin{frame}{One Sample, Two Classifications}
  \small
  \begin{center}
    \begin{tabular}{lrrrr}
      \toprule
      \textbf{Outcome} & \textbf{Email} & \textbf{Search} & \textbf{Social} & \textbf{Total} \\
      \midrule
      Converted & 40 & 55 & 18 & 113 \\
      Not converted & 80 & 70 & 62 & 212 \\
      \midrule
      Total & 120 & 125 & 80 & 325 \\
      \bottomrule
    \end{tabular}
  \end{center}

  \vspace{0.4em}
  \begin{exampleblock}{Independence expected count}
    $E_{ij}=(\text{row total}\times\text{column total})/n$. Converted-email
    expected count: $113\times120/325=41.723$.
  \end{exampleblock}
\end{frame}

\begin{frame}[fragile]{Step 1: Contingency Table}
  \begin{columns}[T]
    \begin{column}{0.42\textwidth}
      \small
      \begin{lstlisting}
E = outer(row, col) / n
      \end{lstlisting}

      \begin{block}{Verified output}
        $n=325$\\
        Converted $E$: $41.723$, $43.462$, $27.815$\\
        All $E\ge 5$: yes
      \end{block}
    \end{column}
    \begin{column}{0.58\textwidth}
      \centering
      \includegraphics[width=\textwidth]{../figures/independence_01_contingency_table.png}
      \vspace{0.2em}
      \scriptsize Observed $2\times 3$ table from Step 1
    \end{column}
  \end{columns}
\end{frame}

\begin{frame}{Independence $\chi^2$ with $df=2$}
  \small
  \[
    \chi^2=\sum\frac{(O_{ij}-E_{ij})^2}{E_{ij}}=10.114985,
    \qquad df=(2-1)(3-1)=2.
  \]
  \[
    p=0.006361,\qquad
    V=\sqrt{\chi^2/(n\min(r-1,c-1))}=0.176417.
  \]
  The test rejects $H_0$: conversion independent of channel. Cramer's V is
  modest.
\end{frame}

\begin{frame}[fragile]{Step 2: Chi-Square Decision}
  \begin{columns}[T]
    \begin{column}{0.40\textwidth}
      \small
      \begin{lstlisting}
chi2, p, df, E = chi2_contingency(table)
V = sqrt(chi2 / (n * (k - 1)))
      \end{lstlisting}

      \begin{block}{Verified output}
        $\chi^2=10.114985$, $df=2$\\
        $p=0.006361$, $V=0.176417$\\
        Decision: reject $H_0$
      \end{block}
    \end{column}
    \begin{column}{0.60\textwidth}
      \centering
      \includegraphics[width=\textwidth]{../figures/independence_02_chi_square.png}
      \vspace{0.2em}
      \scriptsize Cell contributions from Step 2
    \end{column}
  \end{columns}
\end{frame}

\begin{frame}{Step 3: Association Is Not a Causal Claim}
  \begin{columns}[T]
    \begin{column}{0.42\textwidth}
      \small
      Combined $p=0.006361$ (reject independence).

      \vspace{0.3em}
      Desktop stratum $p=0.368546$.

      \vspace{0.3em}
      Mobile stratum $p=0.761939$.

      \vspace{0.5em}
      \textbf{Limit:} channel was not randomly assigned. After stratifying
      by device, the association vanishes. A significant $\chi^2$ is not a
      causal claim.
    \end{column}
    \begin{column}{0.58\textwidth}
      \centering
      \includegraphics[width=\textwidth]{../figures/independence_03_no_causation_limit.png}
      \vspace{0.2em}
      \scriptsize Device confounder from Step 3
    \end{column}
  \end{columns}
\end{frame}

\begin{frame}{Concept Check}
  \begin{enumerate}
    \item Compute $df$ for a $2\times 3$ table.
    \item Why is $E_{11}=113\times120/325$ rather than $np$ from a
      hypothesized mix?
    \item If $p=0.006361$ and $\alpha=0.05$, what is the independence
      decision?
    \item Why do desktop $p=0.368546$ and mobile $p=0.761939$ block a causal
      channel claim?
  \end{enumerate}

  \vspace{0.6em}
  \begin{block}{Connection to Lesson 42}
    The session can continue directly into a chi-square test of homogeneity,
    which uses the same statistic with a different sampling story.
  \end{block}

  \vfill
  \scriptsize Source: OpenStax, \textit{Introductory Business Statistics 2e},
  Chapter 11, Section 11.4. CC BY-NC-SA.
\end{frame}

\appendix



\end{document}
